\documentclass[12pt]{article}
\usepackage{amssymb}
\usepackage[fleqn]{mathtools}
\setlength{\mathindent}{0pt}%stânga

\usepackage{authblk}%pachet pentru introducerea afilierii
\usepackage[]{geometry}
\usepackage{graphicx}%rotește ecuații
\usepackage{amsthm}
\usepackage{amsmath, amssymb}
\usepackage{babel}

\newenvironment{rightdisplay}
  {\begin{flushright}\(\displaystyle}
  {\)\end{flushright}} %ecuație aliniată cu marginea dreaptă

\theoremstyle{definition}
\newtheorem{defi}{Defini\c tia}
\newtheorem{teo}{Teorema}
\newtheorem{prop}{Propozi\c tia}
\newtheorem{lm}{Lema}
\newtheorem{cor}{Corolarul}
\newtheorem{obs}{Observa\c tia}
\newtheorem{ex}{Exerci\c tiul}
\newtheorem{exemp}{Exemplul}
\renewcommand{\proofname}{Demonstra\c tie}
%\renewcommand{\bibname}{Bibliografie}
\renewcommand{\contentsname}{Cuprins}
%\renewcommand{\chaptername}{Capitolul}
\renewcommand{\figurename}{Fig.}

\title{Generalizarea Șirurilor: Fibonacci, Lucas, Pell}
\author{\enspace}
\affil{\enspace}
\date{\enspace}

\begin{document}
	\section*{Șiruri}
		\subsection*{Fibonacci}
		\[\begin{aligned}
			& F_{n+2} = F_{n+1}+F_{n}, F_1=1, F_2=1 \\
			& 1,1,2,3,5,8,... \\
			& \lim\limits_{n\to \infty} \frac{F_{n+1}}{F_n}= \frac{1+\sqrt{5}}{2}
		\end{aligned}\]
		\subsubsection*{Termenul general}
		\[\begin{aligned}			
			& F_{n+2} = F_{n+1} + F_{n}, F_1 = 1, F_2 = 1\\
			& F_{n+2} - F_{n+1} - F_{n} = 0\\
			& \text{Ecuație caracteristică: } x^2 -x -1 = 0\\
			&\Delta = 5 => x_{1,2} = \frac{1 \pm \sqrt{5}}{2}\\
			&F_n = A\left(\frac{1+\sqrt{5}}{2}\right)^n
    				+ B\left(\frac{1-\sqrt{5}}{2}\right)^n\\
    			& F_n=\frac{1}{\sqrt{5}} \left[ \left(\frac{1+\sqrt{5}}{2}\right)^n-\left(\frac{1-\sqrt{5}}{2}\right)^n \right]
		\end{aligned}\]
		\subsection*{Lucas}
		\[\begin{aligned}
			&L_{n+2}=L_{n+1}+L_n, F_1=1, F_2=3\\
			&1,3,4,7,11,...
		\end{aligned}\]
		\subsubsection*{Termenul general}
		\[\begin{aligned}
			&L_{n+2}-L_{n+1}-L_n = 0\\
			&\text{Ecuație caracteristică: }x^2-x-1=0 \\
		&\Delta = 5 \Rightarrow x_{1,2} = \frac{1 \pm \sqrt{5}}{2}
		\Rightarrow L_n = A\left(\frac{1+\sqrt{5}}{2}\right)^n
						+ B\left(\frac{1-\sqrt{5}}{2}\right)^n
		\end{aligned}\]

		\[\begin{aligned}
			&\left.
				\begin{aligned}
					&L_0=A+B=2\\
					&L_n= A\left(\frac{1+\sqrt{5}}{2}\right)
       						+ B\left(\frac{1-\sqrt{5}}{2}\right)
				\end{aligned}
			\right\}
			\Rightarrow
			A\left(\frac{1+\sqrt{5}}{2}\right)+(1-\sqrt{5}) - A\left(\frac{1+\sqrt{5}}{2}\right)=1\\
			&A\sqrt{5}=\sqrt{5} \Rightarrow A=1 \Rightarrow B=1
		\end{aligned}\]
		
		\subsection*{Pell}
		\[\begin{aligned}
			&P_{n+2}=2P_{n+1}+P_n, F_1=1, F_2=2\\
			&1,2,5,12,29,\cdots
		\end{aligned}\]
		\subsubsection*{Termenul general}
		\[\begin{aligned}
			&P_{n+2}-2P_{n+1}-P_n = 0\\
			&\text{Ecuație caracteristică: }x^2-2x-1=0 \\
		&\Delta = 8 \Rightarrow x_{1,2} = \frac{1 \pm 2\sqrt{2}}{2}=1\pm\sqrt{2}
			\ \ \ \begin{aligned}
				&P_2 = 2 \cdot P_1 + P_0\\
				&2 = 2 + P_0 \Rightarrow P_0=0
			\end{aligned}
		\end{aligned}\]

		\[\begin{aligned}
			&\left\{
				\begin{aligned}
					&P_0=A+B=0\\
					&P_n= A(1+\sqrt{2})+B(1+\sqrt{2})=1 
				\end{aligned}
			\right.
			\Rightarrow
			A(1+\sqrt{2})-A(1-\sqrt{2})=1\\
			&2A\sqrt{2}=1 \Rightarrow A=\frac{1}{2\sqrt{2}} \Rightarrow B=-\frac{1}{2\sqrt{2}}
		\end{aligned}\]
		\section*{Forma generală}
			\subsection*{Fibonacci}
				\[\begin{aligned}
				A\ &=\begin{pmatrix}
					1	&	1	\\
					0	&	1	
				\end{pmatrix}\\
				A^2&= \begin{pmatrix}
					1	&	1	\\
					0	&	1	
				\end{pmatrix} \begin{pmatrix}
					1	&	1	\\
					0	&	1	
				\end{pmatrix}=\begin{pmatrix}
					1	&	2	\\
					0	&	1	
				\end{pmatrix}\\
				A^3&=\begin{pmatrix}
					1	&	2	\\
					0	&	1	
				\end{pmatrix} \begin{pmatrix}
					1	&	1	\\
					0	&	1	
				\end{pmatrix}=\begin{pmatrix}
					1	&	3	\\
					0	&	1	
				\end{pmatrix}\\
				&\cdots\\
				A^n&=\begin{pmatrix}
					1	&	n	\\
					0	&	1	
				\end{pmatrix}\\
				\end{aligned}\]
				
				\[\begin{aligned}
				A\ &=\begin{pmatrix}
					1	&	1	\\
					1	&	0	
				\end{pmatrix}\\
				A^2&= \begin{pmatrix}
					1	&	1	\\
					1	&	0	
				\end{pmatrix} \begin{pmatrix}
					1	&	1	\\
					1	&	0	
				\end{pmatrix}=\begin{pmatrix}
					2	&	1	\\
					1	&	1	
				\end{pmatrix}  \\
				A^3&=\begin{pmatrix}
					2	&	1	\\
					1	&	1	
				\end{pmatrix} \begin{pmatrix}
					1	&	1	\\
					1	&	0	
				\end{pmatrix}=\begin{pmatrix}
					3	&	2	\\
					2	&	1	
				\end{pmatrix} \ \ \ \ \ \ \ \ \ \ \ \ \ \ \ \ \ \ \ \ 
				 							A^n = \begin{pmatrix}
													F_{n+1}	& F_{n}\\
													F_n		&	F_{n-1}
											\end{pmatrix}\ \ \ ^{(\boldsymbol{*})}\\
				A^4&=\begin{pmatrix}
					3	&	2	\\
					2	&	1	
				\end{pmatrix} \begin{pmatrix}
					1	&	1	\\
					1	&	0	
				\end{pmatrix}=\begin{pmatrix}
					5	&	3	\\
					3	&	2	
				\end{pmatrix}\\
				A^5&=\begin{pmatrix}
					5	&	3	\\
					3	&	2	
				\end{pmatrix} \begin{pmatrix}
					1	&	1	\\
					1	&	0	
				\end{pmatrix}=\begin{pmatrix}
					8	&	5	\\
					5	&	3	
				\end{pmatrix}\\
				&\cdots\\
				A^n&=\begin{pmatrix}
					F_n + F_{n-1}	&	F_n	\\
					F_n	&	F_{n-1}	
				\end{pmatrix}=\begin{pmatrix}
				F_n & F_n\\
				F_n & 0
				\end{pmatrix}+\begin{pmatrix}
				F_{n-1} & 0\\
				0		& F_{n-1}
				\end{pmatrix} \Rightarrow A^n = F_n \begin{pmatrix}
				1	&	1\\
				1	&	0
				\end{pmatrix}+F_{n-1} \begin{pmatrix}
				1	&	0\\
				0	&	1
				\end{pmatrix}
				\end{aligned}\]
				
	\[\begin{aligned}
	A^n=F_n \cdot A = F_n \cdot A + F_{n-1} \cdot I_2,\  n \ge 1
	\end{aligned}\]
	\[ \text{Din} (\boldsymbol{*}) \Rightarrow  \left\{\begin{aligned} &\det(A^n) = \begin{vmatrix}
	F_{n+1}	& F_{n}\\
	F_{n}	& F_{n=1}
	\end{vmatrix}= F_{n+1} \cdot F_{n-1} - {F_n}^2 = (\det{A})^n \\
	&\det{A} = \begin{vmatrix}
	1&1\\
	1&0
	\end{vmatrix}=-1 \end{aligned}
	\right. \Rightarrow\] \[ \Rightarrow \det{A^n} = F_{n+1} \cdot F_{n-1} - {F_n}^2 = (-1)^n \]
	\begin{obs} 
	\[\begin{aligned}
	\begin{pmatrix}
	F_{n+2}\\
	F_{n+1}
	\end{pmatrix} &= \begin{pmatrix}
	1&1\\
	1&0
	\end{pmatrix} \begin{pmatrix}
	F_{n+1}\\ F_n
\end{pmatrix}\\	
	\begin{pmatrix}
	F_{n+1}\\
	F_{n}
	\end{pmatrix} &= \begin{pmatrix}
	1&1\\
	1&0
	\end{pmatrix} \begin{pmatrix}
	F_{n}\\ F_{n-1}
	\end{pmatrix}\\
	 &\cdots \\
	\begin{pmatrix}
	F_{2}\\
	F_{1}
	\end{pmatrix}\  &= \begin{pmatrix}
	1&1\\
	1&0
	\end{pmatrix} \begin{pmatrix}
	F_{1}\\ F_0
	\end{pmatrix} \xRightarrow{+} \begin{pmatrix}
	F_{n+1}\\
	F_{n}
	\end{pmatrix} = \begin{pmatrix}
	1&1\\
	1&0
	\end{pmatrix}^n \begin{pmatrix}
	F_1\\ F_0 
	\end{pmatrix}=A^n \begin{pmatrix}
	F1 \\ F0
	\end{pmatrix}
	\end{aligned}\]
	\end{obs}
\subsection*{Pell}
\[\begin{aligned}
A \ &=\begin{pmatrix}
2&1\\
1&0
\end{pmatrix}\\
A^2 &= \begin{pmatrix}
2&1\\
1&0
\end{pmatrix}\begin{pmatrix}
2&1\\
1&0
\end{pmatrix}=\begin{pmatrix}
5&2\\
2&1
\end{pmatrix}\\
A^3 &=  \begin{pmatrix}
5&2\\
2&1
\end{pmatrix}\begin{pmatrix}
2&1\\
1&0
\end{pmatrix}=\begin{pmatrix}
12&5\\
5&2
\end{pmatrix}\\
A^3 &= \begin{pmatrix}
12&5\\
5&2
\end{pmatrix}\begin{pmatrix}
2&1\\
1&0
\end{pmatrix}=\begin{pmatrix}
29&12\\
12&5
\end{pmatrix}\\
&\cdots
A^n &= \begin{pmatrix}
P_{n+1}&P_n\\
P_n&P_{n-1}, \ \ n \ge 1
\end{pmatrix} 
\end{aligned}\]
$$\det{A^n} = \begin{vmatrix}
P_{n+1}&P_n\\
p_n & P_{n-1}
\end{vmatrix}= P_{n+1}\cdot P_{n-1} - {P_n}^2 = (-1)^n$$
\[\begin{aligned}
A^n &= \begin{pmatrix}
2P_n + P_{n-1}& P_n\\
P_n&P_{n-1}
\end{pmatrix}=Pn\begin{pmatrix}
2&1\\
1&0
\end{pmatrix}+\begin{pmatrix}
P_{n-1}&0\\
0&P_{n-1}
\end{pmatrix}\\
&= P_n A+P_{n-1}I_2 
\end{aligned}\]
\subsection*{Lucas}





\[
\overset{\scalebox{1.3}{\ \ B}}{
  \overset{\scalebox{2}[1]{\ \rotatebox{90}{\}}}}{
    \begin{pmatrix}
      1 & 2\\
      2 & -1
    \end{pmatrix}
  }
}
\cdot
\overset{\scalebox{1.3}{\ \ A}}{
  \overset{\scalebox{1.6}[1]{\ \rotatebox{90}{\}}}}{
    \begin{pmatrix}
      1 & 1\\
      1 & 0
    \end{pmatrix}
  }
}
=
\begin{pmatrix}
  3 & 1\\
  1 & 2
\end{pmatrix}
\]
%	\maketitle
%	\tableofcontents

%\begin{thebibliography}{3}
%	\bibitem{Angel} \textbf{Angel,A.R., Porter,S.R.}, \textit{A Survey of Mathematics with Applica-
%		tions,}, New York, 1997
%	\bibitem{Batschelet} \textbf{Batschelet,E.,} \textit{Introduction to Mathematics for Life Scientists,} Springer Verlag, Berlin-Heidelberg-New York, 1971.
%	\bibitem{Cret} \textbf{Creț,F.,} \textit{Elemente de Modelare și Matematici Speciale}, Editura MIRTON, Timișoara, 2000
%	
%\end{thebibliography}
%
\end{document}